\scanpage{114}
On a un hom canonique
\[
f_*(L_\bullet)\otimes M_\bullet\ \deleted{\unsure{\otimes}}\longrightarrow f_*(L_\bullet\otimes M_\bullet)
\]
provenant de
\[
M_\bullet-\mathcal Hom(f_*(L_\bullet),f_*(L_\bullet\otimes M_\bullet))
\]
(lui-même provenant de $M-\mathcal Hom(L_\bullet,L_\bullet\otimes M_\bullet))$)
d'où en composant \inserted{avec} \uline{l'hom. can.}
\[
f_*(L_\bullet\otimes M_\bullet)-f_*(C(L_\bullet\otimes M_\bullet))
\]
l'hom voulu. Il faut \uline{vérifier} que
c'est un hom \uline{fonctoriel} pour les catégories
dérivées \unsure{gauches}, \deleted{ce qui est \unsure{le}}

Pour prouver que la flèche est un isomorphisme,
on est ramené au cas où $L_\bullet,M_\bullet$ sont
réduits au degré $0$; \inserted{$Y$ affine,} puis au cas où $M_\bullet$ est
de la forme $i_!(A_{Y'})$, où $i:\unsure{Y'|X'}\longrightarrow\unsure{Y|X}$ est
affine et étale. Dans ce cas d'ailleurs $M$ est
plat et le $\overset{\mathbb L}{\otimes}$ est le $\otimes$ ordinaire. Notons
qu'on a pour tt $P_\bullet$ sur $Y$, et \deleted{\unsure{Q}} $Q_\bullet$ sur $Y$
\[
i_!(Q)\otimes P\simeq i_!(Q\otimes i^*P)
\]
(c'est en\unsure{f.} le cas particulier du th. pour
$i$ étale, qui se vérifie fibre par fibre\ldots) on

\margin{\unsure{Ça} simplifie, en se ramenant au cas où la
base est réduite à un pt, grâce au th. de changt de base
propre (prendre les fibres des deux mbres !)}
\margin{\diagram{\begin{tikzcd}
X'\arrow[r,"j"]\arrow[d,"f'"']&X\arrow[d,"f",dash]\quad L\\
Y'\arrow[r,"i"']&Y
\end{tikzcd}}}
\mspage{113}
